% /Users/pikachu/books/payments-book/src/parts/part6/ch32-change-tracking.tex
\chapter{Как отслеживать изменения и не читать устаревшее}
\label{ch:change-tracking}
\margintoc


\begin{kaobox}[title=Что вы узнаете из этой главы]
  Как выстроить рабочую систему отслеживания изменений в правилах
  и стандартах; какие сигналы приходят из bulletins, release notes
  и регуляторных актов; как жить с restricted-документами и gap между
  public и member-only информацией; и что должно обновляться в книге,
  а что логичнее сразу выносить в online companion.
\end{kaobox}

Найти правильный документ — только половина задачи. В платёжной
индустрии ошибка чаще возникает не потому, что команда вообще
не нашла спецификацию, а потому, что она прочитала правильный
документ не той версии. Отсюда и практический принцип этой главы:
навигация по источникам должна заканчиваться routine обновления,
а не просто списком полезных ссылок.

\section{Сначала определите, какой слой вообще меняется}
\label{sec:ps-tracking-layers}

Не все изменения одинаково опасны. Полезно сразу разделять
четыре слоя.

\begin{itemize}
  \item \textbf{Policy changes}: новые обязанности участников,
        пороги программ мониторинга, retry limits, dispute rules.
  \item \textbf{Protocol changes}: новые поля, message mapping,
        версия 3DS, обновления token or kernel logic.
  \item \textbf{Compliance changes}: sunset dates, новые контролли
        PCI DSS, интерпретации scope.
  \item \textbf{Regulatory changes}: законы, положения, письма
        регулятора, влияющие на допустимость продукта.
\end{itemize}

Эта классификация нужна не ради аккуратности. Она определяет,
кто в команде должен быть owner изменения и как быстро оно
должно попасть в продукт: policy и protocol часто бьют по
конверсии и acceptance, compliance — по аудиту, regulation —
по праву вообще вести бизнес.

\section{Bulletins, release cadence и restricted gap}
\label{sec:ps-tracking-bulletins}

\scheme{Visa} и \scheme{Mastercard} живут по bulletin-driven модели:
сначала публикуется изменение, потом даётся окно на внедрение,
после чего правило становится обязательным. Именно поэтому сеть
нужно читать не только в моменты инцидента, но и по расписанию.

PCI SSC работает похожим образом, но с другой формой сигналов:
standard update, effective date, transition period, supplementary guidance.
EMVCo публикует specification bulletins и новые версии документов.
Банк России добавляет ещё один канал — проекты актов и общественные
обсуждения до вступления нормы в силу.

Самый неприятный сценарий — restricted gap: публичные материалы уже
намекают на изменение, а точные технические детали находятся только
в member-only документе. Здесь рабочий подход обычно такой:
\begin{itemize}
  \item зафиксировать, что изменение существует и когда вступает в силу;
  \item получить owner-документ через участника сети, эквайера
        или процессора;
  \item до получения деталей не придумывать реализацию по пересказам.
\end{itemize}

\begin{kaobox}[title=Важно, colback=red!5, colbacktitle=red!20]
  Restricted gap опасен тем, что команда начинает достраивать
  недостающие технические детали по форумам, PDF-слайдам и словам
  менеджера партнёра. Это почти всегда быстрее в краткосрочной
  перспективе и почти всегда дороже в момент сертификации.
\end{kaobox}

\section{Change triage: что делать с новым правилом}
\label{sec:ps-tracking-triage}

После получения новости о change важно не просто переслать ссылку
в чат, а провести минимальный triage. В платёжной команде полезно
задавать четыре вопроса.

\begin{enumerate}
  \item Это change про policy, protocol, compliance или regulation?
  \item Кого оно затрагивает: issuer, acquirer, gateway, risk,
        reconciliation, legal?
  \item Требует ли оно code change, process change, documentation
        change или только мониторинга?
  \item Какая последняя безопасная дата внедрения?
\end{enumerate}

Из этого triage вырастает нормальный backlog. Без него изменения
обычно падают в одну из двух ловушек: либо их игнорируют до дедлайна,
либо команда начинает паниковать и срочно переделывать не тот слой
системы.

\begin{kaobox}[title=На практике]
  Удобный артефакт change triage — короткая карточка: источник,
  owner-документ, дата вступления в силу, затронутые системы,
  требуемое действие, ссылка на внутреннюю задачу. Такая карточка
  полезнее длинного конспекта, потому что её можно регулярно
  пересматривать на архитектурном или operational review.
\end{kaobox}

\section{Книга против online companion}
\label{sec:ps-tracking-book-vs-companion}

Не каждое изменение должно попадать в печатный текст. Для книги
важно различать материалы с длинным и коротким сроком жизни.

В печатной части лучше жить тому, что редко меняется:
архитурные связи, роли участников, протоколы, инварианты
decision-making. В online companion логично выносить всё,
что быстро стареет:
\begin{itemize}
  \item актуальные пороги VDMP / VFMP / ECM / EFM;
  \item retry limits и effective dates;
  \item payment method mix и benchmark tables;
  \item release-specific fees, thresholds и dates;
  \item полные mappings response / reason codes.
\end{itemize}

Этот принцип полезен не только для книги, но и для внутренних
баз знаний компании. Если материал меняется каждые полгода,
ему не место в длинном narrative-документе без owner-цикла обновления.

\section{Рабочий ритм наблюдения}
\label{sec:ps-tracking-rhythm}

Наблюдение за изменениями лучше работает как ритуал, а не как
реакция на инцидент. Для большинства платёжных команд хватает
простого ритма.

\begin{itemize}
  \item Еженедельно: проверить входящие bulletins, regulatory news
        и PCI/EMVCo updates.
  \item Ежемесячно: провести change triage по новым документам
        и обновить карту рисков.
  \item Ежеквартально: пересмотреть, какие справочные таблицы
        устарели и что нужно вынести или обновить в online companion.
  \item Перед крупным релизом: перечитать owner-документы по тем
        областям, которые релиз реально затрагивает.
\end{itemize}

Такой ритм не делает команду идеальной, но резко уменьшает шанс,
что о новом обязательном поле или пороге она узнает из decline rate,
штрафа или сертификационного отчёта.

\section*{Итоги главы}

\begin{itemize}
  \item Главная проблема в платёжных источниках — не только найти
        документ, но и не перепутать его версию и слой изменения.
  \item Bulletins, release notes, PCI updates, EMVCo bulletins
        и regulatory drafts нужно читать как сигналы к triage,
        а не как информационный шум.
  \item Restricted gap требует дисциплины: сначала получить
        owner-документ через участника сети, потом проектировать
        изменение.
  \item Каждое существенное изменение полезно раскладывать по четырём
        вопросам: что меняется, кого затрагивает, какой нужен action
        и к какому сроку.
  \item Печатная книга должна держать инварианты и архитектуру,
        а быстро стареющие пороги, цифры и даты должны жить
        в online companion.
\end{itemize}
